\documentclass[11pt,a4paper]{article}

\usepackage[utf8]{inputenc}
\usepackage[T1]{fontenc}
\usepackage[english]{babel}
\usepackage{amsmath,amssymb}
\usepackage[top=1.2cm,bottom=1.5cm,left=2cm,right=2cm]{geometry}
\usepackage{graphicx}
\usepackage{booktabs}
\usepackage{hyperref}
\hypersetup{colorlinks=true,linkcolor=black,citecolor=blue,urlcolor=blue}

\title{Question B --- The Bias of the Independence Assumption Between Trees}
\author{Alexandre Boyer}
\date{}

\begin{document}
\maketitle

\noindent The $M$ trees of the forest share the same data stream --- in which
direction does this bias the article's formula, which assumes the $\tau_i$ independent
to compute $\tau_{\mathrm{ARF}} = \min_i \tau_i$? New objects relative to Question A:
$S$ is the complete realization of the stream (one seed); $U_i$ the randomness specific
to tree $i$, identified in Section~\ref{sec:factorization}; $F(s)$ the marginal,
unconditional law of a single tree; $G_S(s) := P(\tau_1>s\mid S)$ its survival function
\emph{given} the stream; $M=10$ the number of trees in the R2 campaign; $\beta$ a
reliability target, the share of drifts one wants the external monitor to catch
($\beta = 0.95$ in the manuscript).

\section{Conditional Factorization}
\label{sec:factorization}

Once $S$ is fixed (the entire trajectory $(x_t,y_t)$ known), what still distinguishes
each tree is $U_i$: the Poisson weights drawn online for bagging (Oza \& Russell, 2001
--- this is not classical bagging by resampling a fixed dataset, which is impossible in
streaming, but a random weight assigned to each new observation), and the random
feature subspaces tested at each node.

The first mechanism is the one the manuscript itself invokes: ``the $M$ trees in the ARF
are trained on the \emph{same} data stream via Poisson-weighted bootstrap (Oza bagging
with $\lambda = 6$)'' (manuscript, Section~III-G, citing Gomes et al.\ 2017). The second
is a defining property of the Adaptive Random Forest in the same
reference\footnote{On the R2 stream it contributes very little, though: the stream
carries two features only ($x_0, x_1$), so a random subspace selects roughly one of two
at each node. Inter-tree diversity there rests almost entirely on the Poisson weights,
which suffices for what follows --- only the i.i.d.\ character of the $U_i$ is needed.
The campaign leaves \texttt{max\_features} at \texttt{river}'s default, which should be
read off the environment that produced the runs rather than assumed.}. These draws are,
by construction, independent across trees: tree $i$ never sees tree $j$'s weights, and a tree newly created at
replacement inherits no information (neither weights nor structure) from the tree it
replaces --- the only link between the two is temporal, the new one having started
learning more recently.

The adaptation time of tree $i$ is a deterministic function of the pair $(S, U_i)$.
Once $S$ is fixed, only the variation in $U_i$ remains, and the $U_i$ are i.i.d.\
across trees. Hence, conditionally on $S$, the $\tau_i$ are indeed independent (and
identically distributed, since all trees share the same procedure).

\begin{center}
\fbox{\parbox{0.93\textwidth}{\textbf{Result.} Defining the conditional survival
function $G_S(s) := P(\tau_1>s \mid S)$, conditional independence gives
\[
  P(\tau_{\mathrm{ARF}}>s \mid S) = P\big(\forall i,\ \tau_i>s \mid S\big) =
  \prod_{i=1}^{M} P(\tau_i>s\mid S) = G_S(s)^M.
\]}}
\end{center}

\section{The Direction of the Bias, via Convexity}
\label{sec:convexity}

We uncondition using the law of total probability:
\[
  P(\tau_{\mathrm{ARF}}>s) = \mathbb{E}_S\big[\,G_S(s)^M\,\big].
\]
The function $x \mapsto x^M$ is convex on $[0,1]$ for $M\ge 1$. Jensen's inequality
gives directly
\[
  \mathbb{E}_S\big[G_S(s)^M\big] \;\ge\; \big(\mathbb{E}_S[G_S(s)]\big)^M = \big(1-F(s)\big)^M,
\]
since $\mathbb{E}_S[G_S(s)] = P(\tau_1>s) = 1-F(s)$ (the marginal law, unconditioned).

\begin{center}
\fbox{\parbox{0.93\textwidth}{\textbf{Result.}
\[
  P(\tau_{\mathrm{ARF}}>s) \;\ge\; \big(1-F(s)\big)^M \quad\Longleftrightarrow\quad
  P(\tau_{\mathrm{ARF}}\le s) \;\le\; 1-\big(1-F(s)\big)^M.
\]
The true probability of having already adapted by time $s$ is always less than or equal
to what the independence formula predicts: the manuscript \textbf{overestimates the
repair speed} --- $\tau_{\mathrm{ARF}}$ is stochastically larger (slower) than the
article claims.}}
\end{center}

\paragraph{Equality condition.} Jensen's inequality is an equality if and only if
$G_S(s)$ is almost surely constant --- that is, if the difficulty of adaptation does
not actually depend on the specific stream realization. Physically, this case
essentially never occurs here: every seed produces a different noise trajectory after
the drift, so $G_S(s)$ genuinely varies across realizations. Independence is therefore
never exact --- only an approximation, systematically optimistic.

\section{Operational Consequence}

Section~\ref{sec:convexity} fixed the \emph{sign} of the error, and two quantities are
affected in opposite directions. The words ``optimistic'' and ``pessimistic'' below refer
to different subjects, which is what makes the pair look contradictory at first sight.

\begin{center}
\begin{tabular}{@{}ll@{}}
\toprule
Quantity & What the independence formula gets wrong \\
\midrule
$\tau_{\mathrm{ARF}}$, the repair speed & it predicts a \emph{faster} repair than the real one \\
$P_{\mathrm{miss}}$, the blind-spot risk & it predicts a \emph{larger} risk than the real one \\
\bottomrule
\end{tabular}
\end{center}

\noindent The two are consistent, and one implies the other: a repair that is in truth
slower leaves the external detector more time to accumulate, so the forest wins the race
less often than the formula claims. Over-stating the repair speed therefore over-states
the blind-spot risk. The manuscript is thus too flattering about the classifier and too
alarming about the architecture --- and our correction, read alone, makes the blind spot
\emph{less} severe than the article reports.

\paragraph{Consequence for $M_{\mathrm{crit}}$.} The manuscript's Corollary~10 turns
Proposition~9 into a prescription on ensemble size, and answers this question: \emph{how
many trees may the forest carry before the external monitor stops being reliable?} It is
evaluated at $\tau^*_{\mathrm{det}} = \lambda/(\Delta e - \delta_P)$, the nominal
accumulation time that Proposition~9 substitutes for the random $\tau_{\mathrm{det}}$ ---
a substitution Question~A examines separately, and which we take as given here.

Requiring
$P_{\mathrm{miss}} = 1 - [1-F(\tau^*_{\mathrm{det}})]^M \le 1-\beta$ gives
$[1-F]^M \ge \beta$, hence $M \ln(1-F) \ge \ln\beta$, and since $\ln(1-F) < 0$ the
division flips the inequality:
\[
  M \;\le\; M_{\mathrm{crit}} := \left\lfloor \frac{\ln\beta}{\ln[1-F(\tau^*_{\mathrm{det}})]} \right\rfloor .
\]
$M_{\mathrm{crit}}$ is therefore a \emph{ceiling} on the number of trees, not a floor ---
which is the physically sensible reading, since more trees mean a faster repair and hence
a worse position for the detector. Adding trees can never help it.

Section~\ref{sec:convexity} established that the true CDF of $\tau_{\mathrm{ARF}}$ lies
everywhere below the independence CDF, so the true $P_{\mathrm{miss}}$ is \emph{smaller}
than the formula's at every $M$. Whatever $M$ the independence certificate admits, the
real forest satisfies the guarantee too.

\begin{center}
\fbox{\parbox{0.93\textwidth}{\textbf{Result.} The corollary is not invalidated but
\textbf{reclassified as a conservative ceiling}: ``at most $M_{\mathrm{crit}}$ trees''
remains a safe prescription, and the true admissible ensemble size is $\ge
M_{\mathrm{crit}}$. It is a valid upper limit, no longer a tight one.}}
\end{center}

\paragraph{One assumption left tacit.} The comparison above holds at a fixed
$\tau^*_{\mathrm{det}}$: it opposes two laws for $\tau_{\mathrm{ARF}}$ against the same
detector behaviour, which is legitimate because Proposition~9 makes
$\tau^*_{\mathrm{det}} = \lambda/(\Delta e - \delta_P)$ a constant. Read as a claim about
the \emph{real} blind-spot risk, the table above assumes more, since the real
$\tau_{\mathrm{det}}$ is computed on the \emph{ensemble's} error stream and therefore
depends on how correlated the trees are. We do not establish that its law is insensitive
to that dependence.

Both effects in fact lower $P_{\mathrm{miss}}$, so the direction of
Section~\ref{sec:convexity} is reinforced rather than threatened. A slower repair leaves
the excess error in place longer, so the monitor accumulates more and alarms earlier; and
correlation weakens bagging's variance reduction, leaving a burstier error stream that
breaches $\lambda$ on noise alone --- the \emph{stochastic resonance} the manuscript
reports for a single tree, which the forest normally suppresses. What weakens is not the
sign but the \emph{meaning} of $P_{\mathrm{miss}}$: part of the decrease is bought with
alarms that win the race without detecting anything. $P_{\mathrm{miss}}$ remains a
well-defined probability, and a poorer proxy for monitoring quality than the formula
suggests.

\paragraph{And the corollary is not rescued either.} On the operative grid the ceiling is
already empty, which our loosening cannot fill. At $\Delta e = 0.33$ the manuscript
measures $\widehat F(\tau^*_{\mathrm{det}}) = 0.49$: a \emph{single} tree has a $49\%$
chance of swapping before the monitor's nominal accumulation time. One tree alone
therefore gives $P_{\mathrm{miss}} = 1 - (1-0.49)^1 = 0.49$, ten times the $5\%$ the
target $\beta = 0.95$ allows, and further trees can only make it worse. No ensemble size
meets the guarantee --- whence $M_{\mathrm{crit}} = 0$, a ceiling of zero trees.

Our correction states that the true ceiling is $\ge M_{\mathrm{crit}}$, which here reads
``$\ge 0$'' and says nothing: rescuing the certificate would require showing the true
ceiling to be $\ge 1$, and a one-sided inequality does not deliver that. The correction
therefore changes the corollary's \emph{status} --- an exact figure becomes a conservative
limit --- without changing its verdict where the manuscript applies it.

\section{The Correlation Trap}

Toy case $M=2$ on $\{1,2\}$: is an anti-monotone coupling compatible with the real
ARF?

\paragraph{The counterexample.} Support $\{1,2\}$, marginals $P(\tau_i=1)=\tfrac12$.
Anti-monotone coupling: $\tau_1=1 \iff \tau_2=2$ (perfectly opposite). Under this
coupling, $\{\tau_1=2\}$ implies $\{\tau_2=1\}$, so
$P(\tau_1{>}1,\ \tau_2{>}1)=P(\tau_1{=}2,\tau_2{=}2)=0$, hence
\[
  \mathrm{Cov}\big(\mathbf{1}\{\tau_1{>}1\},\mathbf{1}\{\tau_2{>}1\}\big) = 0 - \tfrac12\cdot\tfrac12 = -\tfrac14.
\]
Under independence this covariance equals $0$: the anti-monotone coupling thus yields
a strictly negative value.

\paragraph{Why this is impossible in the real ARF.} Reusing the factorization from
Section~\ref{sec:factorization} and the law of total covariance
\[
  \mathrm{Cov}(X,Y) = \mathbb{E}_S[\mathrm{Cov}(X,Y\mid S)] +
  \mathrm{Cov}\big(\mathbb{E}[X\mid S],\mathbb{E}[Y\mid S]\big),
\]
with $X=\mathbf{1}\{\tau_1{>}1\}$, $Y=\mathbf{1}\{\tau_2{>}1\}$: the first term is zero
(conditional independence given $S$, Section~\ref{sec:factorization}). The second is computed directly,
since $\mathbb{E}[X\mid S]=\mathbb{E}[Y\mid S]=G_S(1)$:
\[
  \mathrm{Cov}\big(\mathbf{1}\{\tau_1{>}1\},\mathbf{1}\{\tau_2{>}1\}\big) = \mathrm{Var}\big(G_S(1)\big) \;\ge\; 0 \quad \text{always}.
\]

\begin{center}
\fbox{\parbox{0.93\textwidth}{\textbf{Conclusion --- the trap defused.} The real ARF
can only produce a covariance $\ge 0$ between two trees --- never negative. The
anti-monotone coupling, which requires a covariance of $-1/4$, is therefore
structurally impossible in the real model: the dependence shared through $S$ can only
go one way (independence if $\mathrm{Var}(G_S)=0$, positive dependence otherwise),
never the other. This confirms and locks in the conclusion of Section~\ref{sec:convexity}: no negative
dependence scenario could contradict ``independence is too optimistic about speed''.}}
\end{center}

\end{document}
